\begin{abstract}
Pseudorandom correlation generators and proof systems need binary codes that
combine constant rate and linear distance with fast ordinary and transposed
encoding. They also need predictable data movement. Random linear codes provide
strong distance but require dense computation. Fast structured encoders reduce
this cost, but their sparse dependencies make linear distance difficult to
prove. We ask whether one code family can achieve all of these properties.

We introduce Single-Permutation INterleaved (SPIN) codes. A SPIN encoder
composes a blockwise outer encoder, one global interleaver, and a recursive
inner encoder. Our central result is a Structured SPIN family at native lengths
$N$. Every realized code has rate $1/2$, and its ordinary and transposed
encoders use $O(N)$ bit operations. With probability $1-o(1)$ over setup, its
  minimum distance exceeds $0.11N$. The construction reuses one sampled
  spectrum-certified BA-3 constituent along the outer diagonal, factors the
  interleaver into structured routing stages, and uses the fixed RM2Sub-S19
  inner. The outer interface permits different base constituents and base
  lengths. The current certificate uses extended Golay blocks and the schedule
  $\frac{39}{4}\log_2N+O(1)$. Zero extension gives a family at every
  sufficiently large requested length with rate $1/2-o(1)$, relative distance
  greater than $0.11-o(1)$, and linear work.

The proof uses a finite, class-indexed first-moment framework. The outer stage
supplies class multiplicities, while the interleaver--inner pair supplies
conditional failure envelopes. Two additional families isolate the main proof
mechanisms. At rate $1/2$, Accumulator SPIN achieves relative distance
$0.0037$, and Random SPIN achieves relative distance $0.11002$, under explicit
  logarithmic schedules. A finite-length specialization is under separate
  constituent selection and certification; it is not used by the scalable
  theorem.
\end{abstract}
